% ============================================================
\chapter{Analyse Dimensionnelle et Mise \`a l'\'Echelle}
\label{ch:dimensionnelle}
% ============================================================

En 1883, Osborne Reynolds publie une \'etude sur l'\'ecoulement de fluides
dans des tuyaux. Il d\'ecouvre que le passage du r\'egime laminaire au
r\'egime turbulent ne d\'epend pas s\'epar\'ement de la vitesse, du diam\`etre
ou de la viscosit\'e, mais d'un unique \emph{nombre sans dimension} --- le
nombre de Reynolds. Cette observation illustre le pouvoir de l'\emph{analyse
dimensionnelle}~: en identifiant les combinaisons adimensionnelles des
param\`etres d'un probl\`eme, on r\'eduit drastiquement sa complexit\'e.

% ------------------------------------------------------------
\section{Introduction}
% ------------------------------------------------------------

L'analyse dimensionnelle est un outil puissant qui exploite la
coh\'erence des unit\'es physiques pour simplifier les mod\`eles,
r\'eduire le nombre de param\`etres et d\'ecouvrir des lois d'\'echelle.
Ce chapitre pr\'esente le th\'eor\`eme de Buckingham~$\Pi$, les
techniques d'adimensionnement et leurs applications.

% ------------------------------------------------------------
\section{Dimensions et unit\'es}
% ------------------------------------------------------------

\begin{definition}[Dimension physique]
Toute grandeur physique poss\`ede une \textbf{dimension} exprim\'ee
en fonction des dimensions fondamentales.  Dans le syst\`eme SI, les
dimensions de base sont :
\begin{center}
\begin{tabular}{ccc}
\toprule
Grandeur & Symbole dimensionnel & Unit\'e SI \\
\midrule
Longueur & $L$ & m \\
Masse & $M$ & kg \\
Temps & $T$ & s \\
Temp\'erature & $\Theta$ & K \\
Quantit\'e de mati\`ere & $N$ & mol \\
Intensit\'e \'electrique & $I$ & A \\
\bottomrule
\end{tabular}
\end{center}
\end{definition}

\begin{definition}[Homog\'en\'eit\'e dimensionnelle]
Une \'equation physique est dite \textbf{dimensionnellement homog\`ene}
si chaque terme a la m\^eme dimension.  C'est une condition
n\'ecessaire (mais non suffisante) de validit\'e.
\end{definition}

\begin{exemple}
La vitesse $v$ a pour dimension $[v]=LT^{-1}$.  L'\'energie cin\'etique
$E_c = \frac{1}{2}mv^2$ a pour dimension $[E_c]=ML^2T^{-2}$, qui
est bien celle d'une \'energie.
\end{exemple}

\begin{remarque}
Les arguments de fonctions transcendantes ($\exp$, $\ln$, $\sin$, etc.)
doivent \^etre sans dimension.  Ainsi, dans $e^{-t/\tau}$, le rapport
$t/\tau$ est adimensionnel.
\end{remarque}

% ------------------------------------------------------------
\section{Le th\'eor\`eme de Buckingham $\Pi$}
\label{sec:buckingham}
% ------------------------------------------------------------

\begin{theoreme}[Buckingham $\Pi$]
Soit une relation physique faisant intervenir $n$ grandeurs
$q_1,\ldots,q_n$ dont les dimensions s'expriment \`a l'aide de
$k$ dimensions fondamentales ind\'ependantes.  Alors la relation peut
s'\'ecrire sous la forme d'une \'equation entre $n-k$ groupes
adimensionnels $\Pi_1,\ldots,\Pi_{n-k}$ :
\[
  f(\Pi_1, \Pi_2, \ldots, \Pi_{n-k}) = 0.
\]
\end{theoreme}

\begin{remarque}
Le nombre $k$ est le rang de la \textbf{matrice dimensionnelle},
c'est-\`a-dire la matrice dont les colonnes contiennent les exposants
des dimensions fondamentales pour chaque grandeur.
\end{remarque}

\subsection{M\'ethode pratique}

\begin{enumerate}
  \item Lister toutes les $n$ grandeurs intervenant dans le probl\`eme.
  \item Construire la matrice dimensionnelle et calculer son rang $k$.
  \item Former $n-k$ groupes adimensionnels $\Pi_j$ en combinant les
        grandeurs.
  \item \'Ecrire la loi physique sous forme
        $\Pi_1 = \phi(\Pi_2,\ldots,\Pi_{n-k})$.
\end{enumerate}

\begin{exemple}[P\'eriode d'un pendule simple]
La p\'eriode $\tau$ d\'epend de la longueur $\ell$, de la masse $m$,
de l'acc\'el\'eration gravitationnelle $g$ et de l'amplitude $\theta_0$
(sans dimension).  Les grandeurs dimensionn\'ees sont $\tau$, $\ell$,
$m$, $g$ ($n=4$).  Les dimensions fondamentales sont $M$, $L$, $T$
($k=3$).  Matrice dimensionnelle :
\[
  \begin{pmatrix}
    & \tau & \ell & m & g \\
    M & 0 & 0 & 1 & 0 \\
    L & 0 & 1 & 0 & 1 \\
    T & 1 & 0 & 0 & -2
  \end{pmatrix}
\]
Le rang est $k=3$ (mais la colonne de $m$ est lin\'eairement
ind\'ependante et $m$ appara\^it seule dans $M$, donc $m$ dispara\^it).
On obtient $n-k=1$ groupe :
\[
  \Pi_1 = \tau \sqrt{g/\ell}.
\]
Conclusion : $\tau = C(\theta_0)\sqrt{\ell/g}$ pour une certaine
fonction $C$.  Pour de petites oscillations, $C\approx 2\pi$.
La masse n'intervient pas.
\end{exemple}

% ------------------------------------------------------------
\section{Adimensionnement des \'equations}
\label{sec:adimensionnement}
% ------------------------------------------------------------

\begin{definition}[Adimensionnement]
\textbf{Adimensionner} une \'equation consiste \`a introduire des
variables sans dimension en divisant chaque grandeur par une
\'echelle caract\'eristique :
\[
  \tilde{x} = \frac{x}{x_c}, \qquad
  \tilde{t} = \frac{t}{t_c}, \qquad
  \tilde{u} = \frac{u}{u_c}.
\]
L'\'equation adimensionn\'ee ne fait plus appara\^itre que des
\textbf{nombres sans dimension} qui r\'ev\`elent l'importance relative
des diff\'erents termes.
\end{definition}

\begin{exemple}[\'Equation de l'oscillateur harmonique amorti]
L'\'equation
\[
  m\ddot{x} + c\dot{x} + kx = 0
\]
fait intervenir $m$ (masse), $c$ (amortissement), $k$ (raideur).
Posons $x_c$ une amplitude caract\'eristique et $t_c=\sqrt{m/k}$
(p\'eriode naturelle).  Avec $\tilde{x}=x/x_c$,
$\tilde{t}=t/t_c$ :
\[
  \frac{d^2\tilde{x}}{d\tilde{t}^2}
  + 2\zeta\,\frac{d\tilde{x}}{d\tilde{t}}
  + \tilde{x} = 0,
  \qquad \zeta = \frac{c}{2\sqrt{mk}}.
\]
Le seul param\`etre restant est le \textbf{facteur d'amortissement}
$\zeta$, qui contr\^ole qualitativement le comportement :
\begin{itemize}
  \item $\zeta < 1$ : oscillations amorties (sous-critique),
  \item $\zeta = 1$ : amortissement critique,
  \item $\zeta > 1$ : r\'egime sur-critique (pas d'oscillation).
\end{itemize}
\end{exemple}

\begin{formules}[title=\textbf{Nombres adimensionnels classiques}]
\begin{center}
\begin{tabular}{lcc}
\toprule
Nom & Expression & Signification \\
\midrule
Reynolds & $\mathrm{Re}=\rho v L/\mu$ & inertie / viscosit\'e \\
P\'eclet & $\mathrm{Pe}=vL/D$ & advection / diffusion \\
Mach & $\mathrm{Ma}=v/c_s$ & vitesse / son \\
Froude & $\mathrm{Fr}=v/\sqrt{gL}$ & inertie / gravit\'e \\
Damk\"ohler & $\mathrm{Da}=\tau_{\text{transport}}/\tau_{\text{r\'eaction}}$ & transport / r\'eaction \\
\bottomrule
\end{tabular}
\end{center}
\end{formules}

% ------------------------------------------------------------
\section{Lois d'\'echelle et similitude}
\label{sec:similitude}
% ------------------------------------------------------------

\begin{definition}[Similitude]
Deux syst\`emes physiques sont dits en \textbf{similitude} s'ils
partagent les m\^emes nombres adimensionnels.  L'\'etude d'un
mod\`ele r\'eduit (maquette) peut alors \^etre transpos\'ee au
syst\`eme r\'eel.
\end{definition}

\begin{theoreme}[Invariance d'\'echelle]
Si une grandeur $y$ d\'epend de variables $x_1,\ldots,x_n$ et si
la relation est une loi de puissance :
\[
  y = C\,x_1^{a_1}\,x_2^{a_2}\cdots x_n^{a_n},
\]
alors cette loi est invariante par changement d'unit\'es si et
seulement si les exposants v\'erifient les contraintes dimensionnelles.
\end{theoreme}

\begin{exemple}[Loi de tra\^in\'ee]
La force de tra\^in\'ee $F_D$ sur un objet d\'epend de $\rho$ (densit\'e
du fluide), $v$ (vitesse), $L$ (taille caract\'eristique) et $\mu$
(viscosit\'e).  On a $n=5$, $k=3$ ($M,L,T$), donc $n-k=2$ groupes :
\[
  \Pi_1 = \frac{F_D}{\rho v^2 L^2}, \qquad
  \Pi_2 = \frac{\rho v L}{\mu} = \mathrm{Re}.
\]
Ainsi $C_D = F_D/(\tfrac{1}{2}\rho v^2 A)$ est une fonction de Re
seulement (pour une g\'eom\'etrie donn\'ee).
\end{exemple}

% ------------------------------------------------------------
\section{Application : \'equation de diffusion}
\label{sec:diffusion_dim}
% ------------------------------------------------------------

Consid\'erons l'\'equation de diffusion en une dimension :
\[
  \frac{\partial u}{\partial t} = D\frac{\partial^2 u}{\partial x^2},
  \qquad u(x,0) = u_0(x),
\]
sur un domaine de longueur $L$.  Posons $\tilde{x}=x/L$,
$\tilde{t}=Dt/L^2$, $\tilde{u}=u/U$ (o\`u $U$ est une \'echelle
de concentration).  L'\'equation devient :
\[
  \frac{\partial\tilde{u}}{\partial\tilde{t}}
  = \frac{\partial^2\tilde{u}}{\partial\tilde{x}^2}.
\]
Tous les param\`etres ont disparu : le temps caract\'eristique
de diffusion est $t_c = L^2/D$.

\begin{codeblock}[title=\textbf{Temps de diffusion en Python}]
\begin{minted}{python}
import numpy as np
import matplotlib.pyplot as plt
from scipy.integrate import solve_ivp

# Diffusion 1D par méthode des lignes
L = 1.0    # longueur du domaine
D = 0.01   # coefficient de diffusion
N = 100    # nombre de points
dx = L / (N + 1)
x = np.linspace(dx, L - dx, N)

# Condition initiale : pic gaussien
u0 = np.exp(-((x - 0.5)**2) / (2 * 0.02**2))

def rhs(t, u):
    """Laplacien discret avec conditions de Dirichlet nulles."""
    d2u = np.zeros_like(u)
    d2u[0] = (0 - 2*u[0] + u[1]) / dx**2
    d2u[-1] = (u[-2] - 2*u[-1] + 0) / dx**2
    d2u[1:-1] = (u[:-2] - 2*u[1:-1] + u[2:]) / dx**2
    return D * d2u

sol = solve_ivp(rhs, [0, 5], u0, t_eval=[0, 0.5, 1, 2, 5],
                method='RK45', max_step=0.01)

plt.figure(figsize=(8, 5))
for i, t_val in enumerate(sol.t):
    plt.plot(x, sol.y[:, i], label=f'$t = {t_val}$')
plt.xlabel('$x$'); plt.ylabel('$u(x,t)$')
plt.legend(); plt.title('Diffusion 1D')
plt.tight_layout(); plt.savefig('diffusion_1d.pdf')
\end{minted}
\end{codeblock}

% ------------------------------------------------------------
\section{Perturbations r\'eguli\`eres et ordres de grandeur}
\label{sec:perturbations}
% ------------------------------------------------------------

L'adimensionnement r\'ev\`ele souvent un petit param\`etre $\varepsilon\ll 1$
qui permet une \textbf{approximation perturbative}.

\begin{definition}[D\'eveloppement perturbatif r\'egulier]
On cherche la solution sous la forme
\[
  \tilde{u} = \tilde{u}_0 + \varepsilon\,\tilde{u}_1
             + \varepsilon^2\,\tilde{u}_2 + \cdots
\]
et on identifie les termes ordre par ordre.
\end{definition}

\begin{exemple}
L'oscillateur faiblement non lin\'eaire :
\[
  \ddot{x} + x + \varepsilon x^3 = 0, \qquad 0<\varepsilon\ll 1.
\]
\`A l'ordre z\'ero : $\ddot{x}_0 + x_0 = 0 \Rightarrow x_0 = A\cos t$.
\`A l'ordre un :
\[
  \ddot{x}_1 + x_1 = -x_0^3 = -A^3\cos^3 t
  = -\frac{A^3}{4}(3\cos t + \cos 3t).
\]
Le terme en $\cos t$ est r\'esonnant (terme s\'eculaire).  Cela
signale que le d\'eveloppement na\"if \'echoue \`a long terme ; il
faut utiliser la m\'ethode des \'echelles multiples ou la m\'ethode
de Lindstedt--Poincar\'e.
\end{exemple}

\begin{attention}[title=\textbf{Termes s\'eculaires}]
Un d\'eveloppement perturbatif r\'egulier peut contenir des termes
qui croissent sans borne avec le temps (termes s\'eculaires).  Cela
invalide l'approximation sur des temps longs et n\'ecessite des
techniques sp\'ecialis\'ees.
\end{attention}

% ------------------------------------------------------------
\section{Analyse d'ordres de grandeur}
% ------------------------------------------------------------

\begin{proposition}[Simplification par ordres de grandeur]
Apr\`es adimensionnement, si un terme $\varepsilon\,f(\tilde{x})$
est multipli\'e par $\varepsilon\ll 1$ tandis que les autres termes
sont d'ordre $O(1)$, on peut n\'egliger ce terme en premi\`ere
approximation.
\end{proposition}

\begin{exemple}[\'Ecoulement \`a bas Reynolds]
Pour un fluide newtonien incompressible, les \'equations de
Navier--Stokes adimensionn\'ees contiennent le terme
$\frac{1}{\mathrm{Re}}\nabla^2\mathbf{v}$.  Si $\mathrm{Re}\gg 1$,
la viscosit\'e est n\'egligeable (\'ecoulement d'Euler).  Si
$\mathrm{Re}\ll 1$, l'inertie est n\'egligeable (\'ecoulement de
Stokes) :
\[
  \nabla p = \mu\nabla^2\mathbf{v}.
\]
\end{exemple}

% ------------------------------------------------------------
\section{Exercices}
% ------------------------------------------------------------

\begin{exercice}
En utilisant l'analyse dimensionnelle, d\'eterminer la d\'ependance
de la fr\'equence $f$ d'oscillation d'une goutte sph\'erique en
fonction de sa masse $m$, de son rayon $R$ et de la tension de
surface $\gamma$.
\end{exercice}

\begin{exercice}
Adimensionner l'\'equation de Lotka--Volterra :
\[
  \dot{x} = ax - bxy, \qquad \dot{y} = -cy + dxy,
\]
en posant $\tilde{x}=bx/c$, $\tilde{y}=dy/a$, $\tilde{t}=at$.
V\'erifier qu'on obtient un syst\`eme ne d\'ependant que d'un seul
param\`etre $\alpha = c/a$.
\end{exercice}

\begin{exercice}
La vitesse de s\'edimentation $v_s$ d'une sph\`ere de rayon $R$ et de
masse volumique $\rho_s$ dans un fluide de viscosit\'e $\mu$ et de
masse volumique $\rho_f$ sous gravit\'e $g$ est donn\'ee par la
formule de Stokes.
\begin{enumerate}
  \item Retrouver la d\'ependance $v_s\propto R^2$ par analyse
    dimensionnelle.
  \item Pour quelles valeurs de Re cette formule est-elle valide\,?
\end{enumerate}
\end{exercice}

\begin{exercice}
Consid\'erons l'\'equation de la chaleur avec source :
\[
  \rho c_p \frac{\partial T}{\partial t}
  = \kappa\frac{\partial^2 T}{\partial x^2} + Q_0 e^{-x/\ell}.
\]
\begin{enumerate}
  \item Identifier les \'echelles caract\'eristiques $x_c$, $t_c$, $T_c$.
  \item Adimensionner l'\'equation et identifier les nombres sans dimension.
  \item Discuter les r\'egimes limites.
\end{enumerate}
\end{exercice}

\begin{exercice}
En utilisant le th\'eor\`eme $\Pi$ de Buckingham, montrer que la
perte de charge $\Delta p$ dans un tuyau de longueur $L$ et de
diam\`etre $d$, pour un fluide de masse volumique $\rho$ et de
viscosit\'e $\mu$ s'\'ecoulant \`a la vitesse $v$, s'\'ecrit :
\[
  \frac{\Delta p}{\rho v^2} = \phi\!\left(\mathrm{Re},\,\frac{L}{d}\right).
\]
\end{exercice}
